\section{Escopo e material}

O experimento cujos números este relatório reporta é o run canônico
\url{2026-08-24T09-38-20_5aceb11}, executado no commit \texttt{5aceb11}, de
árvore limpa (\texttt{git\_dirty: false}). Toda tabela, figura e teste aqui pode ser
reproduzido com

\begin{quote}\ttfamily\footnotesize
python scripts/run\_experiment.py --full\\
python scripts/analyze\_results.py results/experiments/<run>\\
python scripts/analyze\_review.py results/experiments/<run>
\end{quote}

\noindent
e conferido contra \texttt{analysis.txt}, \texttt{analysis.json} e
\texttt{review\_analysis.txt} gravados dentro do diretório do run. A análise textual
completa, com o encadeamento dos achados, está em
\path{docs/resultados-oracle-n.md}; este documento é o recorte que responde
diretamente às perguntas do subprojeto.

\paragraph{Procedência.} O run anterior (\texttt{2286fcc}) tinha
\texttt{git\_dirty: true} --- o código real divergia do commit registrado, e nada
publicado a partir dele era reproduzível (ADR 0019). O run canônico substitui-o de
árvore limpa. Antes de qualquer análise, os três modos comuns aos dois desenhos
(\texttt{ga}, \texttt{bagging}, \texttt{rf}) foram comparados coluna a coluna:
\textbf{$3480$ combinações base$\times$fold$\times$modo$\times$métrica de pool
comparadas, $0$ divergências, $\max|\Delta| = 0{,}000000000$}. O run sujo fica
\textbf{validado}: todo número já publicado sobrevive e ganha procedência
reprodutível.

\subsection{Protocolo}

\begin{table}[htbp]
\centering\small
\caption{Protocolo experimental.}
\begin{tabular}{ll}
\toprule
Item & Valor \\
\midrule
Bases & $29$ (das $31$ do catálogo; Ecoli e Glass excluídas pelo portão de viabilidade) \\
Validação & $10$-fold estratificado, \texttt{random\_state=42} \\
Partição interna & treino / validação-\emph{fitness} / DSEL, disjuntas \\
Tamanho do pool & $M=100$ em todos os modos \\
Modos de pool & PGDCS completo, PGDCS-F1/T1, bags aleatórios, Bagging, Random Forest \\
Fusores medidos & $14$ (MVR, combinação suave, $10$ dinâmicos, $2$ estáticos) \\
Métricas & acurácia, precisão e revocação macro, F1 macro, acurácia balanceada \\
Estatística & Friedman + Nemenyi (protocolo Dem\v{s}ar) e Wilcoxon pareado com Holm \\
Unidade pareada & a \emph{base} (folds promediados antes do teste) \\
Total & $1450$ folds, $0$ erros, $16$h$44$ de relógio \\
\bottomrule
\end{tabular}
\end{table}

Duas bases foram excluídas por um portão de viabilidade explícito: a classe mais
rara de Ecoli tem $2$ instâncias e a de Glass tem $9$, ambas abaixo dos $10$ folds,
o que impede a estratificação. A exclusão está registrada no manifesto do run com o
motivo, e não é um descarte silencioso.

\paragraph{Os cinco modos.} Cada um isola uma variável (ADR 0019):

\begin{description}
\item[PGDCS completo.] Perceptrons lineares, algoritmo genético, e a etapa em que o
método \emph{vota} quais medidas de complexidade convêm a cada base
(\texttt{get\_best\_types}) ativa --- a especificação original do método.
\item[PGDCS-F1/T1.] Idêntico, mas com as medidas fixas em F1 e T1 em vez de votadas.
É o modo já publicado em versões anteriores deste relatório; mantido para
comparabilidade e para o critério de aceitação do run acima.
\item[Bags aleatórios.] Mesmos bags estratificados de $50\%$, mesmo Perceptron,
mesma seed, \textbf{sem} busca genética --- a população de geração $0$ do próprio
GA. Isola o efeito da busca.
\item[Bagging.] Árvores de decisão sobre reamostragens \emph{bootstrap}.
\item[Random Forest.] Árvores com subamostragem adicional de atributos.
\end{description}

Os contrastes que os cinco modos destravam --- e cujos resultados aparecem nas
seções seguintes --- são: PGDCS completo \emph{vs} PGDCS-F1/T1 (efeito de votar as
medidas), PGDCS-F1/T1 \emph{vs} bags aleatórios (efeito da busca do GA), bags
aleatórios \emph{vs} Bagging (efeito do classificador-base).

\subsection{Definições usadas}

Para um pool com $M$ classificadores e uma amostra $x$,
\begin{equation}
\On{N}(x) = \mathbf{1}\!\left[\ \sum_{j=1}^{M} \mathbf{1}[\,h_j(x)=y(x)\,] \ \ge\ N \ \right],
\qquad N = 1,\dots,M ,
\end{equation}
de modo que \On{1} é o Oracle tradicional e vale
$\On{1}\ge\On{2}\ge\dots\ge\On{M}$.

Duas quantidades derivadas organizam o relatório:

\begin{description}
\item[Folga.] $\;\On{1}-\mvr$, a parcela do potencial anunciado pelo Oracle que a
votação majoritária não converte.

\item[Folga recuperada.] $\;\dfrac{\text{melhor método dinâmico}-\mvr}{\On{1}-\mvr}$,
a fração dessa folga que um método real alcança. O numerador usa o \emph{máximo por
fold} entre os dez métodos dinâmicos, escolhido no conjunto de teste: é um
\emph{oracle sobre métodos} e portanto uma cota superior generosa do que a seleção
dinâmica entrega.
\end{description}

Uma terceira quantidade mede diversidade de forma comparável entre pools de força
diferente. O Double Fault bruto (\texttt{DF}, a fração de pares que erram juntos)
penaliza pools fracos por construção, porque classificadores que erram mais erram
juntos mais. Normalizando pelo valor esperado sob erros independentes,
\begin{equation}
\dfe, \qquad e = 1 - \text{acurácia individual média},
\end{equation}
obtém-se um índice em que $1$ significa erros independentes e valores altos
significam redundância. O denominador acima assume taxa de erro igual entre
classificadores; a versão exata, $\frac{2}{M(M-1)}\sum_{i<j}e_ie_j$, é usada como
verificação de rigor na seção do objetivo 8.
